\section{Recommendations}
\label{sec:discussion}

For benchmark users and builders, our results suggest six practices. (1) Report \cpr{} (accuracy on items the reader
cannot answer closed-book) next to raw accuracy, together with the reader's mean guess rate. (2) Estimate the guess rate
with the exact evaluation prompt and decoding, not with a separate closed-book prompt. (3) Fix the answer protocol and
report abstention rates; never compare systems whose readers abstain at different rates on raw accuracy alone. (4) Score
unanswerable items by discrimination, never by accuracy averaged into the memory score. (5) Include items whose content
contradicts the prior, ideally as same-slot pairs, and updates in both directions, because a memory's failures are
invisible on facts the model would have guessed anyway. (6) When attributing errors to stages of the memory pipeline,
attribute correct answers as well: a write, update or retrieval failure that a guess rescues is otherwise counted as a
success.

For memory designers, the controlled study locates the loss: whole-memory rewrites drop most items regardless of
congruence, local updates and append-only stores keep them, and a prior-aware write policy (Nemori's predict-calibrate)
kept all items in our setting. Counter-prior facts and updates are exactly what a memory must protect, because they are
the only facts the reader cannot reconstruct; for the same reason they are the only items on which a failure to write,
keep or retrieve a fact is visible.
